% ---- Базовая преамбула: XeLaTeX + кириллица ----
\usepackage{fontspec}
\usepackage{polyglossia}
\setdefaultlanguage{russian}
\setotherlanguage{english}

\setmainfont{Times New Roman}[Ligatures=TeX]
\newfontfamily\cyrillicfont{Times New Roman}[Ligatures=TeX]
\setsansfont{Arial}[Ligatures=TeX]

\usepackage{amsmath,amssymb,amsthm}
\usepackage{geometry}
\geometry{a4paper, top=2cm, bottom=2cm, left=2.5cm, right=1.5cm}
\usepackage{setspace}
\linespread{1.0}\selectfont
\usepackage{booktabs}
\usepackage{graphicx}
\usepackage{enumitem}
\usepackage[hidelinks]{hyperref}
\usepackage{titlesec}
\titlespacing*{\section}{0pt}{2.2ex plus 0.4ex}{1.2ex plus 0.2ex}
\titlespacing*{\subsection}{0pt}{1.8ex plus 0.3ex}{0.8ex plus 0.2ex}
\usepackage{caption}
\captionsetup{font=small}

\setlength{\parindent}{1.25cm}
\usepackage{indentfirst}

% Нумерация формул по разделам
\numberwithin{equation}{section}

% Мягкая вёрстка: убирает выступы строк за поля
\sloppy
\emergencystretch=2em

% Токены действий: Times New Roman не содержит малых капителей,
% поэтому отрисовываем прописными в основном шрифте
\newcommand{\act}[1]{\mbox{\small #1}}

% ---- Графика и таблицы ----
\usepackage{tikz}
\usetikzlibrary{positioning,arrows.meta,fit,backgrounds,calc}
\usepackage{pgfplots}
\pgfplotsset{compat=1.18}
\usepackage{tabularx}
\usepackage{xltabular}
\usepackage{array}
\newcolumntype{Y}{>{\raggedright\arraybackslash}X}
\newcolumntype{C}[1]{>{\centering\arraybackslash}p{#1}}
\newcolumntype{L}[1]{>{\raggedright\arraybackslash}p{#1}}
\usepackage{float}
